% Modernised reading — fonds Alexandre Grothendieck, shelfmark 121, pages 1-72.
%
% This is an INTERPRETATION, not a transcription. Everything the critical
% apparatus carried moves into footnotes. In any dispute of substance the
% transcription governs: whoever cites Grothendieck must cite
% batch-01.fr.tex ... batch-04.fr.tex.
%
% Pass: Opus 5.5 (claude-opus-5-5), 2026-10-03 — first pass, unchecked against the pages by a human.
% The folder was read whole: four batches, pages 1-72, all transcribed.
% Pages with no \page in the transcription: 2, 24, 28, 37, 39, 52, 63, 68
% (blank separators) and 47 (the pink cover « Équisingularité / Cellules »).
%
% Notation fixed for the whole document. F (\mathcal{F}) a family of subsets,
% F_0 its closed members; his conditions keep his letters a) b) c) d) e) f),
% except that the SECOND c) of p. 4 (closure) is called « c bis) » here; the
% conditions alpha beta gamma gamma' epsilon of pp. 6-8 keep his Greek
% letters. The pieces of a partition are « cellules » in §§ 2-4 (pp. 10-20,
% where he writes « feuillets » and « cellules ») and « feuilles » from p. 25
% on (where he writes « feuilles »); « feuille fermée » is the closure of a
% leaf, as he defines it. |X| is the underlying space of a tame object, M the
% tame category, M_n the classes of tame subsets of the cube I^n, I_0 the
% tame interval when it is an object of M; C_n and « C-modèle » for pp.
% 69-72, without the small mark he puts on the C (transcribed C_* there).
%
% Substantive departures from the pages, each footnoted where it occurs:
% — p. 5, condition d): read as « A contains a dense open subset of its
%   closure » (interior relative to the closure); the literal reading
%   (interior in X) would exclude points, which § 2 needs.
% — p. 5, the Prop. (generation by closed members) is stated under a)-e):
%   the page says a)-d) but its proof uses e).
% — p. 6: under alpha, beta, gamma, condition d) is automatic; the page's
%   unfinished delta) is not needed, as its own argument shows.
% — p. 18: « A-bar_i rare dans A-bar_{i+1} » corrected to the reverse.
% — p. 20: « si X compact » read « même si X est compact ».
% — p. 27, Prop. D: false in the topological category (double suspension);
%   footnoted, the reading states it as a question, as the page does.
% — p. 55: the module of tame functions is a module over the tame field only
%   under 3°) (product), a Q-vector space under 1°) 2°).
% — p. 56: the composite's graph needs a projection, which the page omits.
% — p. 61: the « Y x_X Z » the transcription flags as sic is correct, taken
%   over the retractions; no correction.
% — p. 69: the Prop. « X, Y in C_n => X cap Y in C_n » does not follow from
%   (1)-(4) (counterexample in a footnote); stated with stability under
%   inverse images by coordinate maps supplied, which is M'2) of p. 55.
% — p. 69 c): the illegible condition is replaced by what is true (a
%   bijective C-morphism has a C-morphism inverse, by (1)).
% — p. 71: the sheaf of C-morphisms uses the gluing of p. 72, hence (6).
%
% What the folder announces and never establishes: condition f) of p. 23
% (blank); the « Dém. » of p. 48 (broken off); Conjectures 1 and 2 of pp.
% 33-34; Props. A-D of pp. 26-27 (posed as questions); the existence of
% amalgamated sums of tame objects in general (pp. 59-61, 72: reduced to
% conditions 5°) 6°) or to an extension property, not proved from the
% axioms); condition c) of p. 65 for the cube construction of pp. 66-67.

\documentclass[11pt,a4paper]{article}
\input{../preamble/grothendieck}

\folder{121}
\pages{1}{72}
\dating{1974}
\watermark{Édition de démonstration\\interprétation personnelle de l'œuvre}
\foldertitle{Topologie modérée : notes manuscrites (s.d.), lettre (1974) — lecture modernisée du dossier entier}

\begin{document}

\section*{Résumé}

\begin{resume}
Ces feuillets cherchent quelles figures il est raisonnable d'appeler des
« formes » géométriques, et comment les découper proprement en morceaux
simples.

La topologie qu'on apprend d'abord admet des ensembles monstrueux : une courbe
qui remplit un carré, une partie de la droite qui est dense et de complémentaire
dense, une fonction continue qui n'est dérivable nulle part. Ces objets ne
sont pas faux, mais ils ne ressemblent à rien de ce qu'un géomètre dessine. Le
dossier part de l'idée inverse : choisir d'avance une famille de parties
« sages » — celles qu'on obtient à partir de polynômes, de morceaux de droites
et de plans — et ne travailler qu'avec elles. On renonce aux monstres et l'on
gagne en échange des théorèmes de finitude : toute figure sage se découpe en un
nombre fini de cellules, chacune un morceau lisse et d'un seul tenant, collées
le long de leurs bords d'une façon qui ne dépend que de la figure. C'est ce
qu'il appelle la \emph{topologie modérée}.

Le dossier attaque le problème trois fois, par trois portes différentes, et
c'est son intérêt. La première porte (pages 3 à 23) est ensembliste : on se
donne la famille des parties sages comme une algèbre de Boole stable par
adhérence, avec deux exigences — une partie sage a un bord « mince », et on ne
peut pas emboîter indéfiniment des bords dans des bords — et l'on démontre
qu'on peut alors tout découper en cellules, avec un découpage plus grossier que
tous les autres. La deuxième porte (pages 25 à 51) est géométrique : deux
points d'une figure ont « la même singularité » si la figure a, autour de
l'un et de l'autre, la même allure ; regrouper les points de même allure
devrait donner un découpage canonique, et il le conjecture, pour un espace et
même pour une application. La troisième porte (pages 53 à 72) est catégorique :
au lieu de dire quelles parties sont sages, on dit quelles \emph{opérations}
le sont — produits, projections, images — et l'on cherche le petit nombre
d'axiomes qui permettent de reconstruire tout le reste à partir de l'intervalle
$[0,1]$, de ses points, de la moyenne $(x+y)/2$ et du produit $xy$.

Une image aide pour cette dernière porte. Un jeu de construction fournit
quelques pièces de base et quelques façons de les assembler ; les objets qu'on
peut monter sont tout ce qu'on peut monter, rien de plus. Ici les pièces sont
les sous-ensembles sages des cubes $[0,1]^n$, et les règles d'assemblage sont
les produits et les projections. La question est de savoir si l'on peut aussi
\emph{recoller} deux objets le long d'un troisième sans sortir du jeu : c'est la
question des sommes amalgamées, à laquelle les pages 58 à 72 reviennent sans
cesse.

Le dossier est une suite de brouillons, non une rédaction : deux preuves
s'arrêtent net, une condition est annoncée et laissée en blanc, plusieurs
énoncés sont posés comme questions. Il contient aussi ce qui semble être le
brouillon d'une lettre au mathématicien A'Campo, où il annonce une construction conjecturale du
découpage canonique d'une application propre.

Ce que ces pages anticipent a reçu des noms plus tard. Les « structures
modérées » sont très proches de ce qu'on appelle depuis les années 1980 les
\emph{structures o-minimales}, en théorie des modèles ; les découpages en
cellules sont des \emph{stratifications} vérifiant la \emph{condition de
frontière} ; l'équisingularité topologique et le découpage canonique d'une
application sont le domaine de la théorie de Thom–Mather et des théorèmes de
trivialité de Hardt. Dix ans plus tard, l'\emph{Esquisse d'un programme}
(1984) reprend le programme sous le nom de « topologie modérée » ; ce dossier
en est un des chantiers anciens.

\keywords{tame topology, o-minimal structure, definable sets, stratification, frontier condition, cell decomposition, Boolean algebra of sets, finite T0 space, intrinsic stratification, topological equisingularity, Thom-Mather theory, Hardt triviality, double suspension theorem, recollement, manifolds with fibred corners, Lawvere theory, adhesive category, definable Tietze extension, definable triangulation}
\end{resume}

\pagerange{1}{72}
\section*{Le dossier, sa charpente et ses notations}

Soixante-douze pages, dont neuf sans texte suivi : huit feuillets blancs de
séparation (pages 2, 24, 28, 37, 39, 52, 63, 68) et la chemise rose de la page
47, titrée de sa main « Équisingularité / Cellules ». La chemise
grise de la page 1 porte de sa main « Topologie Modérée ». Rien n'est daté ; la
date de 1974 est celle de l'inventaire, qui l'attache à une « lettre ». Le
dossier contient deux brouillons de lettres de sa main : les pages 29 à 36, où
il se réfère deux fois à une lettre précédente à A'Campo et qui ont tout l'air
d'une lettre au même correspondant, et, tête-bêche sur la page 66, le début
biffé d'une lettre de bienvenue en anglais, sans mathématiques, que la
transcription ne reproduit pas.\footnote{La transcription signale que la page
40, un feuillet de cahier, porte en tête le nom « Fujii Guruji Nichidatsu »,
et que la lettre biffée de la page 66 s'adresse au même « Fujii Guruji ».
L'édition n'en tire aucune datation et ne décide pas laquelle de ces pièces est
la « lettre (1974) » de l'inventaire.}

Le fil n'est pas continu : ce sont des moutures successives, sur des papiers et
des encres différents, d'un même projet. On les garde distinctes.

\begin{itemize}
\item \textbf{Pages 1 et 3 à 23 — l'approche ensembliste.} Un ensemble
  $\mathcal{F}$ de parties d'un espace $X$, ses conditions a) à f), les
  décompositions cellulaires, l'existence d'un découpage plus grossier que tous
  les autres, les décompositions régulières, la localisation et le passage aux
  faisceaux. Paragraphes numérotés de sa main, §~1 à §~5.
\item \textbf{Pages 25 à 27 — l'ordre de mobilité.} Une singularité est
  « mobile » si la figure se déplace le long d'une variété en gardant la même
  allure ; quatre propositions posées comme questions.
\item \textbf{Pages 29 à 36 — « Équisingularité », mis au propre.} Partitions
  équisingulières et privilégiées d'un espace et d'un diagramme d'espaces, deux
  conjectures d'existence d'une partition canonique, et la construction
  conjecturale pour une application propre (brouillon de lettre à A'Campo).
\item \textbf{Pages 38 et 40 à 46 — recoller.} Une catégorie de faisceaux
  « recollés » sur un diagramme $Y \leftarrow Z \hookrightarrow X$, puis la
  reconstruction d'un espace stratifié à partir de variétés à bord fibrées les
  unes sur les autres.
\item \textbf{Pages 47 à 51 — « Équisingularité. Cellules ».} Reprise des
  partitions admissibles, du point de vue des feuilles fermées.
\item \textbf{Pages 53 à 67 — la catégorie modérée.} Axiomes M0) à M5) d'une
  catégorie munie d'un foncteur vers les espaces compacts, leur traduction en
  classes $\mathcal{M}_n$ de parties des cubes, la généralisation à une
  catégorie $\mathcal{K}$ quelconque, les sommes amalgamées.
\item \textbf{Pages 69 à 72 — les $C$-modèles.} Une dernière mouture, mise au
  propre et numérotée « 1) », « 2) » : des classes $C_n$ de parties de $R^n$
  et la catégorie qu'elles engendrent.
\end{itemize}

\textbf{Deux homonymies à ne pas confondre.} Le mot « admissible » sert trois
fois : une décomposition cellulaire est $\mathcal{F}$-admissible (page 13) si
ses cellules sont dans $\mathcal{F}$ ; une partition est admissible (page 31,
et page 51 sous une forme un peu différente) si chaque feuille est localement
fermée et d'adhérence réunion de feuilles ; un système de partitions sur un
diagramme est admissible (page 32) si les flèches envoient feuille dans
feuille. Les lettres a) b) c) désignent aussi plusieurs listes de conditions
sans rapport ; on dit chaque fois de quelle liste il s'agit. Enfin le §~1
emploie deux fois la lettre c) : on appelle ici \emph{c~bis)} la seconde
(stabilité par adhérence).\footnote{Le double emploi est sur la page, pages 3
et 4 ; la transcription le garde tel quel.}

\textbf{Vocabulaire.} Une partie d'un espace est \emph{rare} si son adhérence
est d'intérieur vide ; $\overline{A}$, $\mathring{A}$ et
$\partial A = \overline{A} \setminus \mathring{A}$ sont l'adhérence,
l'intérieur et la frontière.\footnote{Il note la frontière tantôt $\dot{A}$
(page 4), tantôt $bA$ (pages 5, 31) ; on écrit partout $\partial A$, sauf aux
pages 40 à 46, où $\partial$ est le bord d'une variété à bord, comme sur la
page.} Les morceaux d'une partition s'appellent \emph{cellules} aux §§~2 à 4
et \emph{feuilles} à partir de la page 25, comme il les appelle.

\pagerange{1}{1}
\section*{La couverture (page 1)}

Sous le titre, des esquisses au crayon : des cubes $I^{n-1} \to I^n \leftarrow
I^{n+1}$, des plongements $X \hookrightarrow I^n$, un espace $X$ présenté comme
quotient d'une somme $\coprod X_i$ par des flèches de recollement, une
relation d'équivalence $R \simeq X \times_Y X$ et, dans un cadre, un
plongement $\beta : X \hookrightarrow M$ soumis à deux conditions : les
composés $\beta \circ p_i$ sont modérés, et $\beta$ est injectif.\footnote{Le
mot du cadre se lit « modulaires » ou « modérées », lecture douteuse, suivi de
deux mots illisibles. On retient « modérées », parce que la condition est
celle que l'axiome M5) de la page 55 imposera : tout objet modéré admet un
morphisme modéré injectif dans un cube. Les esquisses ne sont pas un texte
suivi, et l'on ne s'appuie pas sur elles.} Ce sont les ingrédients de la
troisième mouture du dossier : présenter un objet par recollement de pièces, et
le plonger dans un cube.

\pagerange{3}{9}
\section*{§ 1. Structure modérée sur un espace (pages 3 à 9)}

\subsection*{L'algèbre de Boole et l'anneau des fonctions étagées}

Soit $X$ un ensemble et $\mathcal{F}$ un ensemble de parties de $X$ stable par
réunions finies (a), intersections finies (b) et passage au complémentaire
(c) : une \emph{algèbre de Boole de parties} de $X$. Pour un anneau commutatif
$\Lambda$, soit $\mathcal{A}_{\mathcal{F}}$ l'ensemble des fonctions
$f : X \to \Lambda$ ne prenant qu'un nombre fini de valeurs et dont chaque
fibre $f^{-1}(\lambda)$ est dans $\mathcal{F}$. C'est une sous-$\Lambda$-algèbre
unitaire de l'algèbre $\mathrm{Et}(X, \Lambda)$ des fonctions étagées — les
fibres d'une somme ou d'un produit sont des réunions finies d'intersections
$f^{-1}(\lambda) \cap g^{-1}(\mu)$ — et $\mathcal{F}$ s'en déduit, comme
l'ensemble des $A$ dont la fonction caractéristique est dans
$\mathcal{A}_{\mathcal{F}}$.\footnote{La page énonce ces équivalences sous des
étiquettes a$'$), b$'$), b$''$) qui ne correspondent pas exactement : b$'$)
est invoquée mais n'est écrite que dans une note marginale presque illisible.
On énonce ce qui est vrai. La page écrit aussi « sous-anneau de
$X^{\mathbb{Z}}$ » pour $\mathbb{Z}^X$.}

Quand $\Lambda$ est intègre, de corps des fractions $K$, il affirme que les
algèbres ainsi obtenues sont exactement les traces sur $\mathrm{Et}(X,\Lambda)$
des sous-algèbres de $\mathrm{Et}(X, K)$, c'est-à-dire les sous-algèbres
\emph{saturées} : si $\lambda \neq 0$ et $\lambda \varphi \in \mathcal{A}$, alors
$\varphi \in \mathcal{A}$. C'est juste pour les sous-algèbres unitaires : si
$f \in \mathcal{A}$ prend les valeurs $c_1, \dots, c_k$, l'interpolation de
Lagrange donne $d \cdot \mathbf{1}_{f^{-1}(c_j)} = \prod_{l \neq j} (f - c_l)
\in \mathcal{A}$ avec $d = \prod_{l \neq j}(c_j - c_l) \neq 0$, et la
saturation fait le reste.\footnote{L'hypothèse « unitaire » est ajoutée ; la
page a biffé le mot à la ligne précédente. Elle est nécessaire : sans les
constantes, le produit $\prod (f - c_l)$ n'a pas de sens dans $\mathcal{A}$.}

\subsection*{Les conditions topologiques (pages 4 et 5)}

Supposons maintenant $X$ topologique. On considère :

\begin{itemize}
\item[c~bis)] pour tout $A \in \mathcal{F}$, $\overline{A} \in \mathcal{F}$ ;
  de façon équivalente (par passage au complémentaire), $\mathring{A} \in
  \mathcal{F}$. Alors $\partial A \in \mathcal{F}$, et la partition de $X$ en
  $\mathring{A}$, $\partial A$, $X \setminus \overline{A}$ est formée
  d'éléments de $\mathcal{F}$.
\item[d)] pour tout $A \in \mathcal{F}$, $A$ contient un ouvert dense de son
  adhérence ; autrement dit, $\overline{A}$ privé de l'intérieur de $A$
  \emph{relatif à} $\overline{A}$ est rare dans $\overline{A}$ (et a fortiori
  dans $X$).\footnote{La ligne est très raturée ; « l'intérieur de $A$ est
  dense dans $\overline{A}$ » y est écrit au-dessus d'une première version.
  Lu au pied de la lettre, avec l'intérieur pris dans $X$, l'énoncé
  interdirait à un point d'appartenir à $\mathcal{F}$ dans $\mathbb{R}$, ce
  que le §~2 exige (les cellules minimales sont des fermés, souvent des
  points). L'intérieur est donc relatif à $\overline{A}$, et c'est d'ailleurs
  ainsi que la page 6 l'emploie (« l'intérieur de $C$ dans
  $D = \overline{C}$ est dense dans $D$ »).}
\item[e)] toute suite décroissante $A_1 \supset A_2 \supset \cdots$ de
  fermés appartenant à $\mathcal{F}$, avec $A_{n+1}$ rare dans $A_n$, est
  stationnaire à $\emptyset$.
\end{itemize}

La condition d) est celle qui distingue les parties sages des parties comme
$\mathbb{Q} \subset \mathbb{R}$ ; la condition e) est une dimension finie
déguisée : on ne peut pas emboîter indéfiniment des bords dans des bords.

\textbf{Une remarque sur les générateurs.} Si $\mathcal{L} \subset
\mathfrak{P}(X)$ est stable par réunions et intersections finies et contient
$\emptyset$ et $X$, l'algèbre de Boole engendrée par $\mathcal{L}$ est formée
des réunions finies de parties $A \cap \complement B$, $A, B \in
\mathcal{L}$.\footnote{Les hypothèses de stabilité sur $\mathcal{L}$ sont
ajoutées : sans elles, il faut des intersections de plusieurs $A$ et de
plusieurs $\complement B$. Le renvoi aux conditions, à cet endroit de la page,
est d'une lecture douteuse (« a) à b) »).}

\textbf{Proposition (page 5).} \emph{Si $\mathcal{F}$ satisfait a) à e), elle
est engendrée, comme algèbre de Boole, par l'ensemble $\mathcal{F}_0$ de ses
éléments fermés.}\footnote{La page dit « a) à d) », mais la preuve qu'elle
donne s'arrête grâce à e), et c'est e) qui la fait terminer. On ajoute e) à
l'énoncé.}

Pour $A \in \mathcal{F}$, on écrit $A = \mathring{A} \cup (A \cap \partial A)$.
L'intérieur $\mathring{A} = \complement\, \overline{\complement A}$ est le
complémentaire d'un fermé de $\mathcal{F}$. Reste $A_1 = A \cap \partial A$,
qui est un élément de $\mathcal{F}$ contenu dans le fermé $X_1 = \partial A$, et
$X_1$ est rare dans $X$ (on le montre plus bas). On recommence dans $X_1$ avec
l'algèbre induite $\mathcal{F}|X_1 = \{ B \in \mathcal{F} : B \subset X_1 \}$,
qui satisfait encore a) à e). On obtient une suite de fermés $X \supset X_1
\supset X_2 \supset \cdots$ de $\mathcal{F}$, chacun rare dans le précédent ;
par e) elle s'arrête, et $A$ est une réunion finie de différences de fermés de
$\mathcal{F}$.\footnote{La page écrit $A \cap b\mathring{A}$ là où l'on
attend $A \cap bA$ ; on lit $\partial A$.}

\subsection*{Les fermés suffisent (pages 6 à 8)}

Puisque tout se ramène à $\mathcal{F}_0$, il demande quand un ensemble
$\mathcal{F}_0$ de fermés est exactement l'ensemble des fermés de l'algèbre de
Boole qu'il engendre. Il faut évidemment

\begin{itemize}
\item[$\alpha$)] $\mathcal{F}_0$ stable par réunions finies ;
\item[$\beta$)] $\mathcal{F}_0$ stable par intersections finies ;
\end{itemize}

et il suffit d'y ajouter

\begin{itemize}
\item[$\gamma$)] pour $A, B \in \mathcal{F}_0$, $\overline{A \cap \complement
  B} \in \mathcal{F}_0$,
\end{itemize}

car l'adhérence d'une réunion finie est la réunion des adhérences. Sous
$\alpha$, $\beta$, $\gamma$, \emph{la condition d) est alors automatique} : si
$C = \bigcup_i C_i$ avec $C_i = A_i \cap \complement B_i$, posons
$D_i = \overline{C_i} \in \mathcal{F}_0$ et $E_i = D_i \cap B_i$. Comme $C_i$
est ouvert et dense dans $D_i$, le fermé $E_i = D_i \setminus C_i$ est rare
dans $D_i$, donc dans $D = \bigcup D_i = \overline{C}$, et $E = \bigcup E_i$
est rare dans $D$. Or $D \setminus E \subset C \subset D$ : $C$ contient un
ouvert dense de son adhérence.\footnote{La page amorce une condition
$\delta$), qui devait traduire d) en termes de $\mathcal{F}_0$, et la laisse
inachevée. L'argument qui suit sur la page — celui qu'on vient de donner — montre
qu'elle n'est pas nécessaire. C'est sans doute pourquoi elle est restée en
suspens ; on ne la restitue pas.}

Sous $\alpha$, $\beta$, $\gamma$, l'adhérence d'un $A \in \mathcal{F}$ est le
plus petit élément de $\mathcal{F}_0$ qui le contient : elle se lit dans
$\mathcal{F}_0$, \emph{sans faire appel à la topologie de $X$}. Il en va donc de
même de la rareté : pour $A \subset B$ dans $\mathcal{F}_0$, $A$ est rare dans
$B$ si et seulement si tout élément de $\mathcal{F}_0$ contenant $B \setminus A$
contient $B$. L'axiome e) se récrit :

\begin{itemize}
\item[$\varepsilon$)] toute suite décroissante d'éléments de
  $\mathcal{F}_0$, chacun « $\mathcal{F}_0$-rare » dans le précédent au sens
  qu'on vient de dire, est stationnaire à $\emptyset$.
\end{itemize}

\textbf{Retrouver la topologie.} Inversement, soit $\mathcal{F}_0$ une famille
de parties d'un ensemble nu $X$, stable par réunions et intersections finies et
contenant $\emptyset$ et $X$. Il existe une topologie pour laquelle les
éléments de $\mathcal{F}_0$ sont fermés et vérifient $\gamma$) si et seulement
si

\begin{itemize}
\item[$\gamma'$)] pour $A, B \in \mathcal{F}_0$, il existe un plus petit
  élément de $\mathcal{F}_0$ contenant $A \cap \complement B$.\footnote{La page
  écrit « l'ensemble des $\in \mathcal{F}_0$ qui contiennent $A \cap
  \complement B$ est $\in \mathcal{F}_0$ » ; c'est l'\emph{intersection} de
  ces ensembles qui doit être dans $\mathcal{F}_0$, et la parenthèse qui suit
  sur la page le dit.}
\end{itemize}

On prend pour fermés les intersections quelconques d'éléments de
$\mathcal{F}_0$ — elles sont stables par réunions finies, par distributivité.
C'est la moins fine des topologies qui conviennent, et elle est caractérisée
par la propriété : tout point a un système fondamental de voisinages ouverts
appartenant à $\mathcal{F}$.

\textbf{Le bord est mince (page 8).} Sous a) à d), la frontière $\partial A$
d'un $A \in \mathcal{F}$ est rare dans $X$. Soit en effet $U$ l'intérieur de
$\partial A$, qui est dans $\mathcal{F}$. Si $U$ n'était pas vide, $U \cap A$
serait non vide (car $U \subset \overline{A}$ est ouvert) et d'intérieur vide
(un ouvert contenu dans $U \cap A$ serait dans $\mathring{A}$, disjoint de
$\partial A$) ; mais $\overline{U \cap A} \supset U$, et par d) $U \cap A$
contiendrait un ouvert dense de son adhérence, qui rencontrerait $U$ en un
ouvert de $X$ contenu dans $A$ — contradiction.\footnote{La page mène
l'argument par réduction au cas $\overline{A} = X$, $\mathring{A} =
\emptyset$ ; la fin en est très raturée et ne se lit que par fragments. On
donne une preuve complète du même énoncé.}

\subsection*{La condition f) (page 9)}

\begin{itemize}
\item[f)] pour tout $A \in \mathcal{F}$ localement fermé — c'est-à-dire
  $A = B \setminus C$ avec $B, C \in \mathcal{F}_0$ —, les composantes
  connexes de $A$ sont en nombre fini et appartiennent à $\mathcal{F}$.
\end{itemize}

Elles sont alors ouvertes dans $A$. Sous f), la connexité d'un $A$
localement fermé se lit dans $\mathcal{F}_0$ : $A$ est connexe si et seulement
s'il n'est pas réunion disjointe de deux éléments non vides de $\mathcal{F}$
fermés dans $A$. Et si les ouverts de $\mathcal{F}$ forment une base de la
topologie, f) entraîne que les éléments de $\mathcal{F}$ sont localement
connexes : la composante d'un point dans $U \cap A$, pour $U$ un petit ouvert
de $\mathcal{F}$, est ouverte, étant l'une d'un nombre fini de composantes
fermées.\footnote{Il dit d'abord « du moins ceux qui sont localement fermés »,
puis « en fait cette restriction est inutile ». C'est juste, mais cela utilise
l'extension de f) à tous les éléments de $\mathcal{F}$, qui n'est prouvée
qu'aux pages 18 et 19. Une condition $(\varphi)$, cerclée au-dessus d'une
ligne, semble nommer la décomposition finie en parties
« $\mathcal{F}_0$-connexes » ; lecture douteuse.}

Les conditions a) à f) sont, à la terminologie près, celles qu'on demande
aujourd'hui aux ensembles définissables d'une structure o-minimale sur
$\mathbb{R}$ : algèbre de Boole stable par adhérence, frontière de dimension
plus petite, dimension finie, nombre fini de composantes définissables.
\footnote{Les structures o-minimales viennent de la théorie des modèles des
années 1980 (van den Dries ; Pillay et Steinhorn ; Knight, Pillay et
Steinhorn) ; le livre de van den Dries, \emph{Tame topology and o-minimal
structures} (1998), emprunte son titre à l'\emph{Esquisse d'un programme}. Ces
feuillets, s'ils sont bien de 1974, sont antérieurs à tout cela, et rien n'y
suggère qu'il ait connu les travaux sur les ensembles semi-algébriques ou
sous-analytiques de Łojasiewicz et Hironaka autrement que de réputation : il ne
cite personne. Une différence de fond : ici la topologie est donnée d'avance,
et $\mathcal{F}$ vit dans un seul espace ; une structure o-minimale est une
famille de parties de \emph{tous} les $\mathbb{R}^n$, stable par projection,
ce que seule la troisième mouture du dossier (pages 53 et suivantes) retrouve.}

\pagerange{10}{13}
\section*{§ 2. Décompositions cellulaires (pages 10 à 13)}

Soit $\mathcal{P} = (X_i)_{i \in I}$ une partition finie de $X$, et
$\mathcal{F}_{\mathcal{P}}$ l'algèbre de Boole des réunions de $X_i$ ; les
fonctions de $\mathcal{A}_{\mathcal{F}_{\mathcal{P}}}$ sont celles qui se
factorisent par l'ensemble quotient $X/\mathcal{R} \simeq I$. Pour
$\mathcal{F}_{\mathcal{P}}$, la condition c~bis) équivaut à

\begin{itemize}
\item[(1)] l'adhérence de chaque $X_i$ est une réunion de cellules,
\end{itemize}

qui est la \emph{condition de frontière} des stratifications. On lui adjoint

\begin{itemize}
\item[(2)] deux cellules distinctes ont des adhérences distinctes.
\end{itemize}

Une \emph{décomposition cellulaire} (finie) est une partition finie
satisfaisant (1) et (2) ; elle est \emph{stricte} si de plus

\begin{itemize}
\item[(3)] les cellules sont connexes.
\end{itemize}

\textbf{Les cellules sont localement fermées.} La relation $i \prec j$ si
$\overline{X_i} \subset \overline{X_j}$ — la \emph{relation d'incidence} — est
un préordre, et un ordre grâce à (2). Une cellule minimale $X_{i_0}$ est
fermée : $\overline{X_{i_0}} \setminus X_{i_0}$ est une réunion de cellules
d'adhérence contenue dans $\overline{X_{i_0}}$, donc strictement par (2), ce
que la minimalité interdit. On recommence sur l'ouvert $X \setminus X_{i_0}$
muni de la partition induite, qui satisfait encore (1) et (2). Inversement, (1)
et « les cellules sont localement fermées » entraînent (2) : une cellule est
ouverte et dense dans son adhérence, et deux ouverts denses d'un même espace se
rencontrent.\footnote{La page contient une première tentative, encadrée et
barrée, de dériver (2) d'autre chose, et une addition interlinéaire de deux
lignes ; la remarque « Mais 1° + les $X_i$ loc.\ fermés implique 2° » est dans
la marge.} L'algèbre $\mathcal{F}_{\mathcal{P}}$ est alors engendrée par les
adhérences des cellules, qui vérifient $\alpha$, $\beta$, $\gamma$ ; donc d)
est automatique, par la page 6, et e) l'est aussi puisque $I$ est fini.

\textbf{(3) équivaut à f).} Si f) vaut pour $\mathcal{F}_{\mathcal{P}}$, une
cellule, localement fermée et élément non vide minimal de
$\mathcal{F}_{\mathcal{P}}$, a ses composantes dans $\mathcal{F}_{\mathcal{P}}$,
donc est connexe. Inversement, sous (3), un $A \in \mathcal{F}_{\mathcal{P}}$
localement fermé hérite d'une décomposition cellulaire stricte — (2) passe à
la partition induite, car si $X_i, X_j \subset A$ et $\overline{X_i} \cap A =
\overline{X_j} \cap A$, alors $X_i \subset \overline{X_j}$ et $X_j \subset
\overline{X_i}$ —, et ses composantes connexes sont des réunions de cellules,
puisque les cellules sont connexes.\footnote{Aux pages 11 et 12, il écrit deux
fois « d) » pour f) ; corrigé.}

\textbf{Le quotient fini.} L'espace quotient $I = X/\mathcal{R}$ est un espace
topologique fini. L'adhérence de $\{i\}$ correspond au plus petit fermé saturé
contenant $X_i$, c'est-à-dire à $\overline{X_i}$ par (1) : la relation
d'incidence est la relation de spécialisation de $I$. La condition (2) dit que
$I$ est $T_0$\footnote{Il écrit que $I$ est « primitif » ; lecture douteuse.
La définition qu'il donne, $\overline{\{i\}} = \overline{\{j\}} \Rightarrow
i = j$, est celle d'un espace de Kolmogorov, $T_0$.}, et la condition (1) dit
que l'application quotient $\varphi : X \to I$ est \emph{ouverte}.\footnote{La
page dit que (1) signifie que « $\overline{\varphi^{-1}(J)}$ commute à
l'adhérence » ; on énonce l'équivalence avec l'ouverture de $\varphi$, qui est
ce que la page conclut et qui est exacte : si $U$ est ouvert, le complémentaire
du saturé de $U$ est la réunion des cellules disjointes de $U$, et (1) le rend
fermé ; inversement, si $\varphi$ est ouverte, tout voisinage d'un point de
$X_k \subset \overline{X_i}$ rencontre $X_i$.} D'où l'énoncé, qui clôt le
paragraphe : \emph{les décompositions cellulaires finies de $X$ correspondent
aux quotients finis $T_0$ de $X$ pour lesquels l'application quotient est
ouverte.}

Un espace fini $T_0$ est la même chose qu'un ensemble ordonné fini ; une
application continue $X \to I$ vers un ensemble ordonné muni de la topologie
d'Alexandrov est ce qu'on appelle aujourd'hui un espace stratifié au-dessus de
$I$, et l'ouverture de l'application est précisément la condition de
frontière.\footnote{Ce langage est celui de Lurie (\emph{Higher Algebra},
appendice A) et de la littérature récente sur les espaces stratifiés ; il est
de trente ans postérieur et la page n'en a évidemment rien. Le dictionnaire
entre espaces finis et ensembles préordonnés est d'Alexandrov (1937).}

\pagerange{13}{16}
\section*{§ 3. Décompositions $\mathcal{F}$-admissibles : existence (pages 13 à 16)}

Une décomposition cellulaire est \emph{$\mathcal{F}$-admissible} si ses cellules
sont dans $\mathcal{F}$, c'est-à-dire si $\mathcal{F}_{\mathcal{P}} \subset
\mathcal{F}$.

\textbf{Proposition (page 14).} \emph{Soit $\mathcal{F}$ satisfaisant a) à e), et
$\mathcal{L} \subset \mathcal{F}$ une partie finie. Il existe une décomposition
cellulaire $\mathcal{F}$-admissible dont les éléments de $\mathcal{L}$ sont des
réunions de cellules, et parmi elles une plus grossière que toutes les
autres.}\footnote{Le mot qui remplace « fine », biffé, est peu lisible ; la
transcription offre « grossière » d'après le sens, et c'est bien ce qui est
vrai : il ne peut exister de décomposition plus \emph{fine} que toutes les
autres.}

Soit $X' = \bigcup_{Y \in \mathcal{L}} \partial Y$, fermé rare comme réunion
finie de fermés rares (page 8), et $U = X \setminus X'$. Dans $U$, chaque
$Y \cap U$ est ouvert et fermé ; les atomes $U_i$ de l'algèbre de Boole finie
qu'ils engendrent sont ouverts et fermés dans $U$, donc ouverts dans $X$. On
résout le même problème dans $X'$ pour la famille
$\mathcal{L}' = \{ Y \cap X' \}_{Y \in \mathcal{L}} \cup \{ \overline{U_i} \cap
X' \}_i$ ; la décomposition cherchée de $X$ est formée des cellules de $X'$ et
des $U_i$. Comme $X'$ est rare dans $X$, la récurrence s'arrête par e).

La page dit que le résultat est le plus grossier sans le justifier ; voici
pourquoi c'est vrai. Si $\mathcal{Q}$ est une décomposition admissible rendant
les $Y \in \mathcal{L}$ réunions de cellules, une cellule de $\mathcal{Q}$ qui
rencontre $\partial Y$ y est contenue : si elle est dans $Y$, elle est dans
l'adhérence d'une cellule de $\complement Y$ par (1), donc disjointe de
$\mathring{Y}$, et symétriquement. Donc $\mathcal{Q}$ se restreint à $U$, où
elle raffine les atomes $U_i$, et à $X'$, où elle rend les $\overline{U_i} \cap
X'$ réunions de cellules ; la récurrence conclut.

\emph{Corollaire 1.} Sous a) à e), $\mathcal{F}$ est réunion filtrante des
$\mathcal{F}_{\mathcal{P}}$, $\mathcal{P}$ parcourant les décompositions
$\mathcal{F}$-admissibles. \emph{Corollaire 2.} Ces décompositions, ordonnées
par raffinement, ont des bornes supérieures finies.\footnote{Une condition
« (Dim) », qui n'est autre que e), est d'abord écrite à part puis barrée, et
les énoncés renvoient ensuite à « a) b) c) d) e) ».}

\emph{Remarque (page 16).} Si $\mathcal{F}$ satisfait aussi f), on peut exiger
des cellules connexes : on remplace les $U_i$ par leurs composantes connexes, en
nombre fini et dans $\mathcal{F}$ puisque $U_i$ est un ouvert de $\mathcal{F}$.
Il existe alors une décomposition stricte plus grossière que toutes les autres,
et $\mathcal{F}$ est réunion filtrante des $\mathcal{F}_{\mathcal{P}}$
stricts.

C'est l'énoncé qu'on appelle aujourd'hui l'existence d'une stratification
compatible avec une famille finie d'ensembles définissables, ici sans
condition de régularité et avec une stratification canonique, la plus
grossière.

\pagerange{16}{17}
\section*{§ 4. Décompositions régulières (pages 16 et 17)}

On se donne en plus une partie $\mathcal{F}_r \subset \mathcal{F}$ d'éléments
localement fermés, dits \emph{réguliers} (on pense aux sous-variétés lisses),
avec

\begin{itemize}
\item[(S)] pour tout fermé $Y \in \mathcal{F}$, il y a une plus grande partie
  ouverte de $Y$ appartenant à $\mathcal{F}_r$, notée $Y_{\mathrm{reg}}$, et
  elle est dense : le lieu singulier $Y_{\mathrm{sing}} = Y \setminus
  Y_{\mathrm{reg}}$ est rare dans $Y$.
\end{itemize}

Une décomposition cellulaire est \emph{$\mathcal{F}$-régulière} si, pour tout
fermé $Y$ réunion de cellules, $Y_{\mathrm{sing}}$ est encore une réunion de
cellules ; les cellules sont alors régulières.\footnote{Les trois premières
lignes de la page 17 sont un palimpseste de ratures et d'additions ; on retient
la condition qui s'en dégage et la conséquence qu'il en tire. Celle-ci suppose
implicitement qu'une partie ouverte d'une partie régulière est régulière : une
cellule est ouverte et dense dans son adhérence, donc disjointe du lieu
singulier de celle-ci, donc ouverte dans sa partie régulière.} La proposition
de la page 17 affirme, sous a) à e) et (S), l'existence d'une décomposition
$\mathcal{F}$-régulière rendant les éléments d'une famille finie $\mathcal{L}$
réunions de cellules, et d'une plus grossière ; de même à cellules connexes
sous f).\footnote{La clause « telle que les $Y \in \mathcal{L}$ soient des
réunions de cellules » manque sur la page, mais la preuve s'en sert. Il écrit
aussi « d) » pour f) à deux reprises.} La preuve n'est qu'indiquée : on retire
en outre les lieux singuliers, $X' = \bigcup_{Y} (\partial Y \cup
Y_{\mathrm{sing}})$, et l'on ajoute à $\mathcal{L}'$ les
$(\overline{U_i})_{\mathrm{sing}}$. Le caractère « le plus grossier » n'est
plus argumenté. Les corollaires suivent le modèle du §~3 : $\mathcal{F}$ est
réunion filtrante des $\mathcal{F}_{\mathcal{P}}$ pour $\mathcal{P}$ régulière
(à cellules connexes sous f)), ces décompositions ont des bornes supérieures
finies, et (cas $\mathcal{L} = \emptyset$) $X$ aurait une
\emph{décomposition régulière canonique} — l'analogue, dans ce cadre abstrait,
de la stratification de Whitney canonique d'un ensemble
analytique.\footnote{Rapprochement de l'édition. La stratification de Whitney
minimale d'un ensemble analytique complexe était connue des spécialistes au
début des années 1970 (Whitney, Thom, Mather) ; la page n'en parle pas.}

\pagerange{18}{23}
\section*{§ 5. Localisation et faisceaux (pages 18 à 23)}

\textbf{Restriction à une partie.} Soit $\mathcal{F}$ satisfaisant a) à e) — ou,
« c'est kif-kif », $\mathcal{F}_0$ satisfaisant $\alpha$, $\beta$, $\gamma$,
$\varepsilon$. Pour $A \in \mathcal{F}$, l'algèbre induite
$\mathcal{F}_A = \{ B \in \mathcal{F} : B \subset A \}$, dans l'espace $A$,
satisfait encore a) à e). Pour d) : tout élément de $\mathcal{F}$ est réunion
finie de localement fermés de $\mathcal{F}$, donc $\mathcal{F}_A$ est engendrée
par ses fermés relatifs, et la page 6 s'applique. Pour e) : si $A_1 \supset A_2
\supset \cdots$ sont des fermés relatifs de $A$ avec $A_{i+1}$ rare dans $A_i$,
les adhérences $\overline{A_{i+1}}$ sont rares dans
$\overline{A_i}$\footnote{La page écrit « $\overline{A_i}$ rare dans
$\overline{A_{i+1}}$ » ; c'est l'inverse, et c'est ce que démontre la
parenthèse qui suit sur la page.}, car $A_i \setminus A_{i+1} \subset
\overline{A_i} \setminus \overline{A_{i+1}}$ est dense dans $\overline{A_i}$ ;
e) dans $X$ conclut. Si $\mathcal{F}$ satisfait f), il en va de même de
$\mathcal{F}_A$, et plus fortement \emph{tout} $A \in \mathcal{F}$ a un nombre
fini de composantes connexes, toutes dans $\mathcal{F}$ : $A$ est réunion finie
de localement fermés connexes, et chaque composante est réunion de certains
d'entre eux.

\textbf{Le faisceau des germes.} Supposons que les ouverts de $\mathcal{F}$
forment une base de la topologie. Soit $\underline{\mathcal{F}}$ le
sous-faisceau du faisceau des germes de parties de $X$ engendré par
$\mathcal{F}$, et $\overline{\mathcal{F}} = \Gamma(X, \underline{\mathcal{F}})$
l'ensemble des parties qui sont \emph{localement} dans $\mathcal{F}$ : $X$ est
recouvert par des ouverts $U_i \in \mathcal{F}$ avec $A \cap U_i \in
\mathcal{F}$. On a $\mathcal{F} \subset \overline{\mathcal{F}}$, et
$\overline{\mathcal{F}}$ satisfait a) à d) ; il n'est pas clair qu'elle
satisfasse e), ni qu'elle soit engendrée par ses fermés. Si $X$ est
quasi-compact, $\overline{\mathcal{F}} = \mathcal{F}$ : un recouvrement fini
suffit.

Sans l'hypothèse de base, on définit encore $\underline{\mathcal{F}}$ — ses
sections sur un ouvert $U$ sont les $A \subset U$ qui coïncident localement avec
des éléments de $\mathcal{F}$ — et $\overline{\mathcal{F}}_U = \Gamma(U,
\underline{\mathcal{F}})$. Mais alors, même si $X$ est compact, on n'a plus
$\overline{\mathcal{F}} = \mathcal{F}$\footnote{La page écrit « si $X$
compact, on ne peut affirmer que $\overline{\mathcal{F}} = \mathcal{F}$ » ;
l'exemple qui suit montre qu'il faut lire « même si ».} : pour
$\mathcal{F} = \{\emptyset, X\}$, $\overline{\mathcal{F}}$ est formée des
parties ouvertes et fermées.

\textbf{Le cas d'une décomposition cellulaire (pages 20 et 21).} Prenons
$\mathcal{F} = \mathcal{F}_{\mathcal{P}}$. Une partie ouverte et fermée d'une
cellule fermée est dans $\overline{\mathcal{F}}$ ; pour avoir
$\overline{\mathcal{F}}_U = \mathcal{F}_U$ pour tout ouvert $U \in \mathcal{F}$,
il faut donc que toutes les cellules soient connexes. Cela suffit : un $U \in
\mathcal{F}$ est une réunion de cellules, et si $A \subset U$ est localement une
réunion de cellules, $A \cap X_i$ est ouvert et fermé dans $X_i$ pour chaque
cellule $X_i$, donc égal à $X_i$ ou vide.\footnote{La page 21 écrit « une
partie à la fois ouverte et fermée de $A$ » ; c'est de $X_i$ qu'il s'agit.}
Pour un ouvert $U$ \emph{quelconque}, les $A \in \overline{\mathcal{F}}_U$ sont
des réunions de composantes connexes des $X_i \cap U$, mais pas toutes ces
réunions : sur $X = \mathbb{R}$, avec les cellules $\{0\}$ et $\mathbb{R}^{*}$,
la demi-droite $]0, +\infty[$ — composante de $\mathbb{R}^{*}$ — n'est pas dans
$\overline{\mathcal{F}}$, car près de $0$ elle ne coïncide avec aucun élément
de $\mathcal{F}$. Et avec des cellules connexes : sur le cercle
$\mathbb{R} \cup \{\infty\}$, avec les cellules $\{0\}$ et son complémentaire,
et $U = \mathbb{R}$, on retrouve le même phénomène.\footnote{La page note la
demi-droite $\mathbb{R}^{+}$ et le cercle $\widehat{\mathbb{R}}$, puis un
symbole surchargé lu comme $\mathbb{S}^1$. « n'est pas » est d'une lecture
douteuse, mais c'est ce que l'exemple exige.}

\textbf{La condition f) n'est pas locale (page 22).} Elle passe aux ouverts de
$\mathcal{F}$ mais pas aux ouverts quelconques. Si $A \in
\overline{\mathcal{F}}_U$ est localement fermé dans $U$, il est localement
connexe, ses composantes sont ouvertes, et localement finies dans $A$ si $A$
est localement compact ; il demande si elles sont localement finies dans $U$,
et laisse la question.

\textbf{Le point de vue des faisceaux.} Il part alors d'un sous-faisceau
$\underline{\mathcal{F}}$ du faisceau des germes de parties, satisfaisant a) à
d) et une forme locale de e) : pour tout ouvert $U$ et toute suite décroissante
de fermés de $U$, sections de $\underline{\mathcal{F}}$, chacun rare dans le
précédent, la suite est localement finie — tout quasi-compact $K \subset U$ ne
rencontre qu'un nombre fini de termes. Alors $\Gamma(X, \underline{\mathcal{F}})$
satisfait a) à d), et e) si $X$ est quasi-compact. Plus commode est l'axiome

\begin{itemize}
\item[e$_n$)] il existe $n$ tel que toute telle suite $X_0 \supset \cdots
  \supset X_{n+1}$ ait $X_{n+1} = \emptyset$ — $X$ est « de
  $\underline{\mathcal{F}}$-dimension $\leqslant n$ » —,
\end{itemize}

qui passe à tous les $\mathcal{F}_A$, $A$ fermé, et entraîne e).

Le faisceau $\underline{\mathcal{F}}$ peut être strictement plus gros que celui
qu'engendrent ses sections globales. Exemple : les germes de parties
analytiques — on entend : de l'algèbre de Boole qu'elles engendrent — d'un tore
complexe de dimension 2 sans fonctions méromorphes non
constantes.\footnote{Justification de l'édition : un tel tore ne contient
aucune courbe. Une courbe sur un tore de dimension 2 est de carré positif ou
nul ; de carré strictement positif, le tore serait projectif, de carré nul, il
fibrerait sur une courbe elliptique, et dans les deux cas il aurait des
fonctions méromorphes non constantes. Ses
seules parties analytiques fermées sont donc les ensembles finis de points et
le tore entier, alors que localement il a autant de germes de courbes qu'on
veut.} Une condition f) est ensuite annoncée (« Nous envisageons maintenant une
condition : f) ») et laissée en blanc ; la page s'arrête là, et l'approche
ensembliste avec elle.

\pagerange{25}{27}
\section*{L'ordre de mobilité d'une singularité (pages 25 à 27)}

Une nouvelle mouture commence, qui ne prend plus une famille de parties pour
donnée, mais cherche à la \emph{définir} à partir de la géométrie.

Soit $X$ un espace topologique et $Y \subset X$ une variété topologique (« sous-espace
non singulier »). On dit que $X$ est \emph{équisingulier le long de $Y$} si
tout $y \in Y$ a un voisinage ouvert $U$ dans $X$ homéomorphe à un produit
$N \times Y_U$ ($Y_U = Y \cap U$, $(N, a)$ un espace pointé), l'homéomorphisme
envoyant $Y_U$ sur $\{a\} \times Y_U$ de façon compatible à l'identité de
$Y_U$. En mots : le long de $Y$, la figure a partout la même allure, et cette
allure est celle du germe $(N, a)$.

Un point $x$ est \emph{$\nu$-mobile} s'il existe une telle $Y \ni x$ avec
$\dim_x Y \geqslant \nu$ ; son \emph{ordre de mobilité} $\nu(x)$ est le plus
grand tel $\nu$, qui existe si $X$ est localement de dimension finie. Les
points d'ordre $0$ sont les singularités \emph{immobiles}, ou
éliminables.\footnote{« ou immobile » est écrit au-dessus d'un mot biffé,
lecture incertaine du mot biffé.} L'ensemble $\mathrm{Mob}^i(X)$ des points
$i$-mobiles est \emph{ouvert} — dans une carte $N \times Y_U$, chaque point
$(n, y)$ est $i$-mobile le long de $\{n\} \times Y_U$ —, d'où une suite
décroissante d'ouverts $X = \mathrm{Mob}^0 \supset \mathrm{Mob}^1 \supset
\cdots$ et la suite croissante des fermés $F_i(X) = X \setminus
\mathrm{Mob}^{i+1}(X)$.\footnote{La page écrit $F_0 \subset F_1 \subset F_3$ ;
on lit $F_2$.}

Il se demande, pour un $X$ donné, si les énoncés suivants sont valables (on
pose $F_{-1} = \emptyset$ et $Y_i = F_i \setminus F_{i-1}$, l'ensemble des
points d'ordre exactement $i$) :

\begin{itemize}
\item[A.] $Y_i$ est une variété de dimension $i$ en chacun de ses points, et
  $X$ est équisingulier le long de $Y_i$.
\item[B.] Localement près de chaque point de $Y_i$, $(X, Y_i) \simeq
  Y_i \times (N, a)$ avec $a$ immobile dans $N$, et le type du germe $(N, a)$ ne
  dépend que de la composante connexe de $Y_i$.
\item[C.] Le germe de $Y_i$ en $y$ est le plus grand germe de variété le long
  duquel $X$ est équisingulier ; les autres sont les germes de sous-variétés de
  $Y_i$.\footnote{« plongées » et « autres » sont d'une lecture douteuse.}
\item[D.] Si $Y$ est une variété de dimension $i$ et $N$ un espace, $(y, n)$
  est $\nu$-mobile dans $Y \times N$ si et seulement si $n$ est $(\nu -
  i)$-mobile dans $N$ ; donc $Y_\nu(Y \times N) = Y \times Y_{\nu-i}(N)$.
\end{itemize}

Ce qu'il construit là est, dans le cadre linéaire par morceaux, la
\emph{dimension intrinsèque} d'un point d'un polyèdre et la
\emph{stratification intrinsèque} qu'elle définit ; dans ce cadre les énoncés
A à D sont du type de ceux que la théorie des polyèdres établit.\footnote{La
stratification intrinsèque des polyèdres figure dans la littérature linéaire
par morceaux de la fin des années 1960 (Armstrong, Akin) ; l'édition n'a pas
vérifié que les énoncés A à D s'y trouvent sous cette forme exacte, et ne
l'affirme pas.} Dans le cadre \emph{topologique}, où la page se place, D est
faux : si $\Sigma$ est une sphère d'homologie de dimension 3 non simplement
connexe, le sommet du cône $c\Sigma$ est immobile, mais $c\Sigma \times
\mathbb{R}^2$ est homéomorphe à $\mathbb{R}^6$, où tout point est
$6$-mobile — c'est le théorème de la double suspension de Cannon et Edwards.
\footnote{Ce théorème est postérieur au dossier (Edwards, milieu des années
1970 ; Cannon, 1979), et la page ne pouvait pas le connaître. Deux notes en
marge, près de B et de D, portent un mot lu « divisible » sous réserve
(« $N$ divisible $\Leftrightarrow$ $Y \times N$ divisible ») : c'est
peut-être précisément la question de pouvoir simplifier un facteur, qu'il
pressentait délicate. Il pose d'ailleurs A à D comme des questions
(« la validité des propositions “suivantes” »), non comme des théorèmes, et la
page 29 quitte le cadre topologique pour le cadre linéaire par morceaux.}

\pagerange{29}{36}
\section*{Équisingularité, mise au propre (pages 29 à 36)}

\subsection*{Le contexte et la définition (page 29)}

Les « espaces » sont maintenant les espaces linéaires par morceaux — ou
algébriques, ou analytiques —, « variété » et « morphisme » s'entendant dans
le contexte choisi. Pour $Y \subset X$ localement fermé, $X$ est équisingulier
le long de $Y$ si tout $y \in Y$ a un voisinage ouvert $X'$ et un espace pointé
$(Z, a)$ avec un carré commutatif

\begin{tikzcd}
Y' = X' \cap Y \arrow[r, hook] \arrow[d, "\simeq"'] & X' \arrow[d, "\simeq"] \\
Y' \times \lbrace a \rbrace \arrow[r, hook] & Y' \times Z
\end{tikzcd}

Plus généralement, pour un diagramme d'espaces $\mathfrak{X} = (X_i)_{i \in I}$
indexé par un ensemble ordonné, et $Y$ localement fermé dans $X_{i_0}$, on
demande que le diagramme formé des $X_j$ pour $j \leqslant i_0$ soit localement
le long de $Y$ le produit par $Y$ d'un diagramme $(Z_j)$ d'espaces, $Z_{i_0}$
pointé, par des plongements ouverts $Z_j \times Y' \to X_j$.\footnote{Il écrit
une fois $Z'_j$ et une fois $Z'_{i_0}$ pour $Z_j$, $Z_{i_0}$ ; on unifie.}

\subsection*{Filtrations et partitions (pages 30 à 32)}

Une filtration croissante $(X_i)_{i \geqslant 0}$ par des fermés, de réunion
$X$, est \emph{équisingulière} si, pour chaque $i$, le diagramme
$X_0 \to X_1 \to \cdots \to X$ est équisingulier le long de $V_i = X_i
\setminus X_{i-1}$. Les \emph{feuilles} sont les composantes connexes des
$V_i$ ; elles partitionnent $X$, et (proposition de la page 30) l'adhérence
d'une feuille est une réunion de feuilles — finie si $X$ est compact et les
$X_i$ localement connexes. Les adhérences s'appellent \emph{feuilles fermées} ;
$F \mapsto \overline{F}$ est bijective ; l'inclusion des feuilles fermées
définit la relation d'incidence ; et $\overline{F} \setminus F$, réunion des
feuilles incidentes à $F$ et distinctes d'elle, est le \emph{bord} de
$F$.\footnote{Il dit aussi « feuilles ouvertes » pour les feuilles
restreintes, ouvertes dans leur adhérence et non dans $X$.} Une seconde
proposition (page 31) caractérise les filtrations équisingulières : il faut et
il suffit que l'adhérence de toute feuille soit une réunion de feuilles, et que
le diagramme formé de $X$ et de toutes les feuilles fermées, avec les
inclusions, soit équisingulier le long de chaque feuille.\footnote{Le mot
« inclusions » est écrit au-dessus d'un mot biffé, lecture douteuse ; de même
« équisingulier » dans b).} Il en vient à
privilégier la partition sur la filtration (page 32) : plusieurs filtrations
donnent la même partition, et dans le contexte linéaire par morceaux il y a au
moins quatre façons canoniques de filtrer une partition donnée (par dimension,
par minimalité itérée, et les deux duales).

Une partition $(F_i)$ de $X$ est \emph{admissible} si chaque $F_i$ est
différence de deux fermés et chaque $\overline{F_i}$ est réunion de feuilles —
c'est la décomposition cellulaire du §~2, sans finitude — et
\emph{équisingulière} si de plus le diagramme formé de $X$, des
$\overline{F_i}$ et de leurs inclusions est équisingulier le long de chaque
feuille. Elle est \emph{privilégiée} si elle est équisingulière et si les
feuilles sont des variétés sans bord.\footnote{La définition de « privilégié »
est dans une note marginale verticale, de lecture fragmentaire ; « admissible »
est aussi dans la marge. Le sens de « privilégié » qu'on retient est celui
qu'imposent les deux conjectures.} Sur un diagramme d'espaces, un système de
partitions est admissible si chaque flèche envoie feuille dans feuille ; il
définit alors un diagramme d'ensembles ordonnés de feuilles, et il est
équisingulier si le diagramme des applications entre feuilles fermées l'est le
long de chaque feuille.

\subsection*{Deux conjectures (pages 33 et 34)}

Dans le contexte « par morceaux » :

\textbf{Conjecture 1.} \emph{Tout espace $X$ admet une partition privilégiée, et
parmi elles une plus grossière que toutes les autres, dite canonique.}

Cela ne suffit pas : il faudra des partitions privilégiées plus fines qu'une
partition donnée, ou rendant des fermés donnés réunions de feuilles, ou
compatibles avec une application. Tout cela rentre dans :

\textbf{Conjecture 2.} \emph{Soit $I$ un ensemble ordonné fini et
$(X_i)_{i \in I}$ un diagramme d'espaces dont les flèches sont propres. Il
existe des partitions privilégiées de ce diagramme, et parmi elles une moins
fine que toutes les autres, dite canonique.}

\emph{Remarque.} Pour un système de partitions équisingulier et une flèche
propre $f : X_i \to X_j$, l'image d'une feuille $F$ est la feuille $\Phi$
entière qui la contient : elle est ouverte dans $\Phi$ par équisingularité,
fermée par propreté, et $\Phi$ est connexe. De plus $F \to \Phi$ est une
fibration localement triviale, et même $f^{-1}(\Phi)$ avec sa partition
induite ; comme $F$ et $\Phi$ sont des variétés, les fibres en sont
aussi.\footnote{La fin de la page 34 est très raturée, et l'argument de
fermeture n'est donné que par fragments ; on en retient la structure, qu'il
conclut sur la page 35.}

\subsection*{Une construction conjecturale pour une application propre (pages 35 et 36)}

« Dans la précédente lettre à A'Campo, j'ai donné une description
conjecturale de la “partition privilégiée canonique” d'un espace $X$, en
termes de points “immobiles”. »\footnote{Cette lettre précédente n'est pas
dans le dossier. Le correspondant est vraisemblablement Norbert A'Campo, alors
spécialiste des singularités ; la page écrit « A-Campo ». La description par les
points immobiles est celle des pages 25 à 27.} Voici, pour une application
propre $f : X \to Y$, $Y$ compact, celle qu'il propose :

\begin{itemize}
\item[1.] Soit $U$ le plus grand ouvert de $Y$ qui soit non singulier et
  au-dessus duquel $f$ soit une fibration localement triviale. Il faut prouver
  que $Y' = Y \setminus U$ est un sous-espace rare, donc de dimension plus
  petite. Les composantes connexes de $U$ sont les feuilles ouvertes de $Y$.
\item[2.] Au-dessus de $U$, les feuilles de $X$ sont celles de la partition
  privilégiée canonique de $V = f^{-1}(U)$, localement le produit de $U$ par
  la partition canonique d'une fibre.
\item[3.] Si $F_i$ sont ces feuilles, les traces $\overline{F_i} \cap
  f^{-1}(Y')$ doivent être des réunions de feuilles. On est ramené à
  $f^{-1}(Y') \to Y'$, de dimension plus petite, mais muni en plus de ces
  fermés et de leurs intersections finies : la récurrence demande donc
  d'emblée l'énoncé pour un espace muni d'une famille finie de fermés.
\end{itemize}

Il note que ce cas plus général « peut probablement se faire en paraphrasant
celle proposée dans la 1re lettre ». L'étape 1 est un théorème de
trivialité générique : dans les cadres semi-algébrique et sous-analytique,
c'est le théorème de Hardt (1980), et dans le cadre o-minimal sa
généralisation.\footnote{Rapprochement de l'édition. La partie existence de la
conjecture 2 est de ce type ; l'existence d'une partition \emph{la plus
grossière} pour une application est une autre affaire, que ces théorèmes ne
donnent pas, et l'édition ne prétend pas qu'elle soit établie.}

\pagerange{38}{38}
\section*{Recollement de faisceaux (page 38)}

Une page de formules sans texte, sur le diagramme $\mathfrak{X} = (Y
\xleftarrow{p} Z \overset{i}{\hookrightarrow} X)$. Un faisceau sur
$\mathfrak{X}$ est un triple $(\mathcal{F}, \mathcal{G}, \alpha)$ formé d'un
faisceau $\mathcal{F}$ sur $Y$, d'un faisceau $\mathcal{G}$ sur $X$ et d'un
morphisme $\alpha : p^{*}\mathcal{F} \to i^{*}\mathcal{G}$. Il note les foncteurs

\[
  i'^{*}(\mathcal{F}, \mathcal{G}, \alpha) = \mathcal{F}, \qquad
  i'_{*}(\mathcal{F}) = (\mathcal{F}, 0, 0), \qquad
  p'^{*}(\mathcal{F}, \mathcal{G}, \alpha) = \mathcal{G}, \qquad
  p'_{*}(\mathcal{G}) = (p_{*} i^{*} \mathcal{G}, \mathcal{G}, \mathrm{can}),
\]

et $q_{!}(\mathcal{G}) = (0, \mathcal{G}, 0)$ ; $q^{!}$ reste inachevé. Ces
formules sont cohérentes : $i'_*$ est adjoint à droite de $i'^*$, $p'_*$ adjoint
à droite de $p'^*$ (un morphisme vers $p'_*\mathcal{G}$ est déterminé par sa
composante sur $X$, par adjonction $p^* \dashv p_*$), et $(0, \mathcal{G}, 0)$
est adjoint à gauche de $p'^*$.\footnote{Vérification de l'édition ; la page ne
dit rien de ces adjonctions. Le nom $q_!$ suggère une extension par zéro depuis
$X \setminus Z$, que le schéma de la page dessine ($q : X \setminus Z \to
\mathfrak{X}$) ; c'est alors à un faisceau sur $X \setminus Z$, prolongé par
zéro, qu'il faut l'appliquer.} C'est la catégorie de recollement de SGA~4,
exposé~IV, transportée à un diagramme ; elle décrit les faisceaux sur le
cylindre de l'application $p$ recollé à $X$ le long de $Z$. En bas de page :
les strates $X_{I,J} = X_I \setminus X_J$ d'un espace stratifié, les quatre
opérations $\alpha_*$, $\alpha^*$, $\alpha_!$, $\alpha^!$ entre elles et le
complexe dualisant — le formalisme des six opérations sur un espace stratifié,
simplement nommé.

\pagerange{40}{46}
\section*{Reconstruire un espace stratifié (pages 40 à 46)}

Ces pages, surtout faites de croquis, cherchent à décrire un espace stratifié
par des données lisses : chaque strate, privée d'un voisinage tubulaire
ouvert de ses strates inférieures, est une variété à bord, et son bord est
fibré au-dessus des strates inférieures.

\textbf{La donnée (page 40).} $I$ est un ensemble ordonné fini. Pour chaque
$i \in I$, $V_i^{\circ}$ est une variété à bord (sans bord si $i$ est
minimal). Pour $j < i$, une partie $\partial_j V_i^{\circ}$ du bord est une
sous-variété à bord, fibrée sur $V_j^{\circ}$, de fibres des variétés à bord
(sans bord si $j$ est un prédécesseur immédiat de $i$), et $\partial
V_i^{\circ}$ est réunion des $\partial_j V_i^{\circ}$ pour $j$ prédécesseur
de $i$. Les bords se décomposent à leur tour : $\partial(\partial_j V_i^\circ)$
contient les $\partial_j V_i^\circ | \partial_k V_j^\circ = \partial_{j,k}
V_i^\circ$, fibrés sur $V_k^\circ$, et ainsi de suite le long des chaînes
$j > k > l > \cdots$.\footnote{Les indices $\partial_{i,j}$,
$\partial_{\emptyset,j}$ et l'indice du « fibré des bords » sont d'une lecture
incertaine ou illisibles ; on ne garde que la structure par chaînes, que les
croquis de la page et ceux de la page 46 (une suite $V_i^{\circ} \leftarrow
V_{ij}^{\circ} \leftarrow V_{ijk}^{\circ} \leftarrow \cdots$ avec des carrés
cartésiens) confirment.} C'est ce qu'on appelle aujourd'hui une variété à
coins fibrés, ou une structure de fibration itérée.\footnote{Le nom vient de la
littérature des années 2000 sur les espaces stratifiés (Albin, Leichtnam,
Mazzeo, Piazza, à la suite de Melrose) ; les « données de contrôle » de Mather
(1970) et la « résolution » des espaces de Thom–Mather par Verona (années 1980)
en sont des formes plus anciennes. La page n'en cite aucune.} La page 40 est
un feuillet de cahier titré « Fujii Guruji Nichidatsu » ; les pages 41, 42, 44
et 45 sont presque entièrement des dessins.

\textbf{Trois strates (page 43).} Pour une tour $U \subset V \subset W$ de
dimensions $r$, $r+s$, $r+s+t$, avec $U$, $V \setminus U$, $W \setminus V$
réguliers et l'équisingularité le long des strates, la donnée est :

\begin{itemize}
\item[1°)] $U$, régulier de dimension $r$ ;
\item[2°)] $U'$ fibré sur $U$ à fibres régulières de dimension $s-1$ ;
\item[3°)] $V^{\circ}$ variété à bord, de bord $U'$ ;
\item[4°)] $V^{\circ\prime}$ fibré sur $V^{\circ}$ à fibres régulières de
  dimension $t-1$, dont le bord $V^{\circ\prime}|U'$ est fibré sur $U$ ;
\item[5°)] $U''$ autre fibré sur $U$, à fibres des variétés à bord de dimension
  $s+t-1$ ;
\item[6°)] un isomorphisme de fibrés sur $U$ entre $\partial U''$ et $\partial
  V^{\circ\prime}$, d'où par recollement l'espace $V' = V^{\circ\prime}
  \sqcup_\partial U''$ ;
\item[7°)] $W^{\circ}$ variété à bord, de bord $V'$.
\end{itemize}

On reconstruit $V = V^{\circ} \sqcup_{U'} U$ (on écrase le bord $U'$ sur $U$
par la fibration) et $W = W^{\circ} \sqcup_{V'} V$.\footnote{Les dimensions de
5°) sont lues sur un bloc serré, exposants incertains. Le diagramme de la page
est redessiné dans la transcription d'après une esquisse dont toutes les
flèches ne se lisent pas ; on ne le reprend pas.} Les pages 44 à 46 étendent
cela à une tour $U_0 \subset U_1 \subset \cdots \subset U_n$, avec les bords
$\partial_{n-1} U_n^{\circ}, \dots, \partial_0 U_n^{\circ}$ fibrés sur
$U_{n-1}^{\circ}, \dots, U_0^{\circ}$, et la page 45 porte le mot
d'ordre : « PL variétés à bord fibrées sur variétés à bord ».

C'est le « dévissage » des structures stratifiées que l'\emph{Esquisse d'un
programme} décrira comme l'un des théorèmes à attendre d'une topologie
modérée : l'équivalence entre espaces stratifiés modérés et systèmes de
variétés à bord recollées.

\pagerange{48}{51}
\section*{Équisingularité. Cellules (pages 47 à 51)}

Sous une chemise rose titrée de sa main, une reprise.

\textbf{Contextes (page 48).} Quatre sont nommés : L.P.M., SA.P.M.,
$\mathrm{Alg}_{\mathbb{C}}$, $\mathrm{An}_{\mathbb{C}}$.\footnote{Les deux
premiers sigles ne sont pas développés. L.P.M.\ est vraisemblablement
« linéaire par morceaux », SA.P.M.\ peut-être « semi-algébrique » ou
« semi-analytique par morceaux » ; l'édition ne tranche pas.} L'énoncé est
celui du §~3, localisé : pour une famille localement finie de sous-espaces
fermés $T_\alpha \subset X$, il existe une partition admissible \emph{stricte}
(à feuilles connexes) dont les $T_\alpha$ sont réunions de feuilles, une plus
grossière que toutes les autres, et sa formation commute à la restriction aux
ouverts. La démonstration, par récurrence sur $\dim X$ — trivial si $\dim X
\leqslant 0$ ; ensuite on retire $X^1 = \bigcup_\alpha \partial T_\alpha$ —,
s'interrompt sur « Soit ».\footnote{Les lignes de la démonstration sont
surchargées d'additions et de passages biffés, et la lecture en est très
incertaine ; on n'en donne que le plan, qui est celui de la page 14.}

\textbf{L'ordre et les ouverts (page 49).} Pour une partition admissible et un
ouvert $U$, soit $I_U$ l'ensemble des feuilles qui rencontrent $U$. L'ordre
d'incidence de $I_U$ calculé dans $U$ est celui de $I$ : si $F_i$ rencontre
$U$ et $\overline{F_i} \cap U \subset \overline{F_j} \cap U$, un point de
$F_i \cap U$ est dans $\overline{F_j}$, donc $F_i \subset \overline{F_j}$ par la
condition de frontière.\footnote{La page passe par un argument plus long,
local au voisinage d'un point, et conclut « Oui, car $F_k$ est ouvert dans
$\overline{F_i}$ ! ». Elle écrit aussi « induite par $U$ » pour « induite par
$I$ ».} En conséquence, si $F_i$ rencontre $U$, $i$ est maximal dans $I$ si et
seulement s'il l'est dans $I_U$ ; et pour une partition admissible localement
finie, \emph{une feuille est ouverte si et seulement si elle est maximale} :
si $F_i$ est maximale, $X \setminus F_i$ est la réunion localement finie des
fermés $\overline{F_j}$, $j \neq i$, qui ne rencontrent pas $F_i$.\footnote{La
page 50 aboutit à cette équivalence sur une ligne serrée qui finit par un
point d'interrogation ; elle est vraie, sous l'hypothèse de finitude locale.}

\textbf{Stabilité par restriction (page 50).} Il distingue trois conditions sur
une partition admissible : (a) feuilles connexes, (b) feuilles non singulières,
(c) équisingularité. Les deux dernières, comme l'admissibilité, passent à un
ouvert, mais pas (a). La bonne condition stable est

\begin{itemize}
\item[(C)] pour tout ouvert $U$, la famille des composantes connexes de tous
  les $U \cap F_i$ est une partition admissible de $U$,
\end{itemize}

c'est-à-dire : ces composantes sont localement finies, et si une composante
$F_{i\alpha}$ de $U \cap F_i$ a une adhérence qui rencontre une composante
$F_{j\beta}$, elle la contient.\footnote{L'addition interlinéaire qui donne la
finitude locale est de place incertaine dans la phrase. Une filtration (III)
de la page 49, encadrée, donne une forme filtrée de la même condition ; ses
indices sont peu lisibles et l'on ne la reprend pas.}

\textbf{Cellules (page 51).} Une partition $(F_i)_{i \in I}$ d'un espace $X$
en \emph{feuilles} est admissible si (I) :

\begin{itemize}
\item[(i)] elle est localement finie ;
\item[(ii)] chaque $\overline{F_i}$ est une réunion de feuilles ;
\item[(iii)] chaque $F_i$ est ouvert dans $\overline{F_i}$.
\end{itemize}

Deux feuilles distinctes ont des adhérences distinctes, et la partition se
lit sur la famille des feuilles fermées $\Phi_i = \overline{F_i}$, qui est
caractérisée par (II) :

\begin{itemize}
\item[(0)] les $\Phi_i$ recouvrent $X$ ;
\item[(i)] la famille est localement finie ;
\item[(ii)] les $\Phi_i$ sont fermés ;
\item[(iii)] $\Phi_i \cap \Phi_j$ est une réunion de $\Phi_k$ ;
\item[(iv)] si $\Phi_j \subsetneq \Phi_i$, $\Phi_j$ est rare dans $\Phi_i$.
\end{itemize}

Réciproquement, (II) suffit, pour une famille de fermés non vides deux à deux
distincts\footnote{Hypothèse ajoutée ; sans elle, $F_i$ peut être vide ou deux
indices donner la même feuille.} : on pose $F_i = \Phi_i \setminus
\bigcup_{\Phi_j \subsetneq \Phi_i} \Phi_j$. Tout $x$ est dans un unique $F_i$
— celui du plus petit $\Phi_i$ contenant $x$, qui existe par (i) et (iii) —,
$F_i$ est localement fermé, la famille est localement finie, $\overline{F_i}
= \Phi_i$ par (iv) (une réunion localement finie de fermés rares est rare), et
$\Phi_i$ est réunion de feuilles : si $F_j$ rencontre $\Phi_i$ en $x$, le plus
petit fermé de la famille contenant $x$ est $\Phi_j$, donc $\Phi_j \subset
\Phi_i$ par (iii).

\pagerange{53}{57}
\section*{La catégorie modérée (pages 53 à 57)}

Troisième mouture : la structure modérée n'est plus une famille de parties
d'un espace donné, mais une catégorie d'objets avec leur espace sous-jacent.

\subsection*{Axiomes (page 53)}

On se donne une catégorie $\mathcal{M}$ et un foncteur $X \mapsto |X|$ de
$\mathcal{M}$ vers les espaces compacts, avec

\begin{itemize}
\item[M0)] le foncteur est « transportable » : on peut transporter la structure
  d'un objet le long d'un homéomorphisme ;
\item[M1)] les limites projectives finies existent dans $\mathcal{M}$ et
  $|\ |$ les respecte ;
\item[M2)] $|\ |$ est conservatif (donc fidèle).
\end{itemize}

Un morphisme $f$ est alors un monomorphisme si et seulement si $|f|$ est
injectif, et les sous-objets de $X$ s'identifient à certaines parties de
$|X|$, les parties \emph{modérées}, stables par intersections finies. Une
application $|X| \to |Y|$ est modérée si elle provient d'un morphisme ; c'est le
cas si et seulement si son graphe est une partie modérée de $|X \times Y|$. On
demande encore :

\begin{itemize}
\item[M3)] l'image d'une partie modérée par une application modérée est
  modérée.
\end{itemize}

Tout morphisme se factorise alors de façon unique en $X \xrightarrow{f'} Y'
\xrightarrow{i} Y$ avec $|f'|$ surjectif et $i$ un monomorphisme, et $f'$ est
un épimorphisme \emph{effectif}. La clé est un lemme : si $|f'|$ est surjectif
et $H \circ |f'|$ est modéré, alors $H$ est modéré, car le graphe de $H$ est
l'image du graphe de $H \circ |f'|$ par $|f'| \times \mathrm{id}$, et M3)
s'applique.\footnote{La continuité de $H$, qu'il ne mentionne pas, vient de ce
qu'une surjection continue entre compacts séparés est une application
quotient. La page écrit « $h : Y' \to |Z|$ », mélangeant objet et espace ; on
lit $h : Y' \to Z$.}

\subsection*{Les ind-objets et la droite réelle (pages 53 et 54)}

Un ind-objet \emph{spécial} est une suite de monomorphismes $X_1 \to X_2 \to
\cdots$ telle que chaque $|X_i|$ soit dans l'intérieur d'un $|X_j|$ dans tous
les $|X_k|$ suivants. Sa limite $\varinjlim |X_i|$ est un espace localement
compact dénombrable à l'infini, dont les $|X_i|$ forment un système
fondamental de compacts. On en tire la notion de structure
\emph{pseudo-modérée} sur un espace localement compact, et l'image directe
d'une partie pseudo-modérée par un morphisme pseudo-modéré \emph{propre} est
pseudo-modérée.

On demande alors une structure pseudo-modérée sur $\mathbb{R}$ pour laquelle
les $I_n = [-n, n]$ sont modérés et la multiplication par un entier non nul est
un isomorphisme. Comme $\mathbb{R} = \text{``}\varinjlim\text{''}\, I_n$ et que
$I_n \simeq I_1$ par $t \mapsto t/n$, les transitions devenant la multiplication
par $n/(n+1)$, cela revient à une structure modérée sur $I_1 = [-1, 1]$ pour
laquelle les homothéties $t \mapsto rt$, $r \in \mathbb{Q}^*$, $|r| \leqslant 1$,
sont modérées. On suppose de plus

\begin{itemize}
\item[M4)] 1°) le point $1$ est modéré ; 2°) la moyenne $(x, y) \mapsto
  (x+y)/2$, $I_1 \times I_1 \to I_1$, est modérée ; 3°) le produit $(x, y)
  \mapsto xy$ est modéré.
\end{itemize}

Sous 1°) et 2°), les points rationnels et les translations rationnelles sont
modérés, et les applications modérées d'un objet dans $\mathbb{R}$ forment un
$\mathbb{Q}$-espace vectoriel. Sous 3°) en plus, les points modérés de
$\mathbb{R}$ forment un corps $\mathbb{R}_{\mathrm{mod}}$ — l'inverse de $a$
est la projection de la partie modérée $\{ (a, y) : ay = 1 \}$ —, et ces
applications forment une $\mathbb{R}_{\mathrm{mod}}$-algèbre.\footnote{La page
55 dit que, « moyennant 1°, 2° », elles forment un module « sur le corps
$\mathbb{R}_{\mathrm{mod}}$ » : la multiplication par un scalaire modéré
demande le produit, c'est-à-dire 3°). On corrige l'attribution des hypothèses ;
« sous-module » et « sous-anneau algèbre » sont d'ailleurs de lecture
douteuse. La dernière ligne de la page 54, coupée par le bord du feuillet,
annonce que les intervalles à extrémités modérées sont modérés.} Enfin :

\begin{itemize}
\item[M5)] tout objet admet un morphisme modéré injectif dans un cube $I^n$.
\end{itemize}

\subsection*{Tout se lit dans les cubes (pages 55 à 57)}

Sous M5), $\mathcal{M}$ se reconstruit à partir des ensembles $\mathcal{M}_n$
des parties compactes modérées de $I^n$, $I = [0, 1]$, qui vérifient

\begin{itemize}
\item[M$'$0)] les éléments de $\mathcal{M}_n$ sont compacts ;
\item[M$'$1)] $\mathcal{M}_n$ est stable par intersections finies ;
\item[M$'$2)] stabilité par images réciproques par les applications de
  coordonnées $I^n \to I^m$ ;
\item[M$'$3)] stabilité par images directes par ces applications.
\end{itemize}

Les « applications simpliciales » de la page sont les applications
$I^J \to I^{J'}$ induites par les applications d'ensembles finis
$J' \to J$ : projections, diagonales, permutations de
coordonnées.\footnote{Il dit que $\mathcal{M}_*$ est un « ensemble
cosimplicial » ; c'est plutôt un foncteur sur la catégorie des ensembles finis,
puisque toutes les applications d'ensembles finis interviennent, pas seulement
les croissantes.} Une structure modérée sur un compact est une classe de
plongements dans des cubes à image dans les $\mathcal{M}_n$, deux plongements
étant équivalents si le graphe de l'application de transition est dans
$\mathcal{M}_{n+n'}$ ; une application est modérée si son graphe, lu dans les
cubes, l'est.

\textbf{La catégorie $\mathcal{M}_0$ (pages 56 et 57).} Partant de classes
$\mathcal{M}_n$ vérifiant M$'$1) à M$'$3), on définit $\mathcal{M}_0$ : ses
objets sont les couples $(n, K)$, $K \in \mathcal{M}_n$, et ses morphismes
$(n,K) \to (n',K')$ les applications $f : K \to K'$ de graphe dans
$\mathcal{M}_{n+n'}$ — ce qui force $f$ continue, son graphe étant compact. Le
graphe de l'identité est l'image de $K$ par la diagonale (M$'$3)). Pour
$f : K \to K'$ et $g : K' \to K''$, le graphe de $g \circ f$ est la
projection sur $I^{n+n''}$ de $p^{-1}(\Gamma_f) \cap q^{-1}(\Gamma_g) \subset
I^{n+n'+n''}$, donc est modéré par M$'$1), M$'$2), M$'$3).\footnote{La page
écrit $\Gamma_{fg} = p^{-1}(\Gamma_f) \cap q^{-1}(\Gamma_g)$, qui est une
partie de $I^{n+n'+n''}$ et non le graphe ; la projection sur $I^{n+n''}$
manque, et M$'$3) est nécessaire en plus de M$'$1) et M$'$2), qu'il cite seuls.}
Le foncteur $(n, K) \mapsto K$ est fidèle et conservatif ; il y a un objet final
$I^0$, des produits fibrés et des noyaux, calculés comme dans les ensembles ;
d'où M1) et M2), M3) par M$'$3), et M5) par construction. Enfin M4) se
traduit :

\begin{itemize}
\item[M$'$4)] pour $q$ rationnel, $|q| \leqslant 1$, le graphe de $x \mapsto
  qx$ est dans $\mathcal{M}_2$ ; $\{1\} \in \mathcal{M}_1$ ; les graphes de la
  moyenne et du produit sont dans $\mathcal{M}_3$.\footnote{Avec $I = [0,1]$,
  comme à la page 55, seuls les $q \geqslant 0$ ont un sens : la condition est
  écrite pour l'intervalle $[-1,1]$ de la page 54.}
\end{itemize}

Ce que M$'$0) à M$'$4) décrivent est, à la compacité près, une
\emph{structure} sur $\mathbb{R}$ au sens de la théorie des modèles, contenant
les graphes de l'addition et de la multiplication : une famille de parties de
tous les $\mathbb{R}^n$ stable par intersections, produits et projections.
L'axiome d'o-minimalité — les parties modérées de la droite sont les
réunions finies de points et d'intervalles — n'est pas ici : il est remplacé,
dans la première mouture, par les conditions de finitude d), e), f).

\pagerange{58}{62}
\section*{Changer de catégorie de base ; sommes amalgamées (pages 58 à 62)}

\textbf{(A) Une base quelconque (page 58).} On remplace les espaces compacts par
une catégorie $\mathcal{K}$ à limites projectives finies, et l'on se donne un
foncteur $\varphi : \mathcal{M} \to \mathcal{K}$ exact à gauche, conservatif,
et un objet $I_0$ de $\mathcal{M}$ dans une puissance duquel tout objet se
plonge. À équivalence près, cela revient à se donner un objet $I$ de
$\mathcal{K}$ et des ensembles $\mathcal{M}_n$ de sous-objets de $I^n$ stables
par intersections finies, par images réciproques d'applications de
coordonnées, et par image d'une application de coordonnées qui reste un
monomorphisme sur $K$.\footnote{Sa lettre pour la catégorie de base est un K
simple ; on écrit $\mathcal{K}$, comme la transcription, pour la distinguer des
parties $K$.}

\textbf{(B) Images.} Si tout morphisme de $\mathcal{K}$ se factorise en un
épimorphisme effectif suivi d'un monomorphisme, et si les épimorphismes
effectifs sont universels, on peut demander la stabilité par images
quelconques (3°~bis), qui est M3).

\textbf{(C) Réunions.} Si les sous-objets ont des bornes supérieures finies, on
peut demander que les $\mathcal{M}_n$ y soient stables et qu'il existe un
monomorphisme modéré $e \sqcup e \to I$ ; cela revient à ce que $\mathcal{M}$
ait des sommes disjointes finies respectées par $\varphi$.

\textbf{(D) Sommes amalgamées (pages 59 à 61).} Un morphisme $f : Y \to X$ est
un épimorphisme si et seulement si la codiagonale $X \sqcup_Y X \to X$ est un
isomorphisme ; si $\mathcal{M}$ a les mêmes sommes amalgamées que $\mathcal{K}$
et que $\varphi$ les respecte, $\varphi$ reflète donc les épimorphismes. La
question devient : étant donnés des monomorphismes modérés $Y \leftarrow X
\rightarrow Z$, trouver sur la somme amalgamée $S = Y \sqcup_X Z$ calculée dans
$\mathcal{K}$ une structure modérée qui en fasse une somme amalgamée dans
$\mathcal{M}$.\footnote{À la page 59, il écrit « $X \xrightarrow{f} X$ » ; la
transcription lit $Y \to X$ d'après la suite, et l'on fait de même.}

Il suppose, dans $\mathcal{K}$, que les sommes amalgamées le long de
monomorphismes existent, que leurs flèches structurales sont des monomorphismes
et que le carré est cartésien, et il veut de plus un carré cartésien de
monomorphismes $X \to Y, Z \to S'$ dont $S$ soit la réunion de $Y$ et de $Z$.
Il en déduit que tout monomorphisme de $\mathcal{K}$ est effectif — c'est le
noyau de la paire $Y \rightrightarrows Y \sqcup_X Y$ — et que les sommes
amalgamées sont universelles. Ce sont presque mot pour mot les propriétés
des \emph{catégories adhésives}.\footnote{Notion de Lack et Sobociński (2004) :
une catégorie où les sommes amalgamées le long de monomorphismes existent et
sont des carrés de Van Kampen ; on y montre que ces carrés sont cartésiens et
que les monomorphismes sont réguliers. Rapprochement de l'édition.} Pour la
structure modérée de $S$, il ajoute deux conditions sur l'intervalle :

\begin{itemize}
\item[5°)] $I_0$ est injectif dans $\mathcal{M}$ : toute application modérée
  d'un sous-objet vers $I_0$ se prolonge ;
\item[6°)] tout sous-objet $Y \hookrightarrow X$ est l'image réciproque d'un
  point $\sigma_n : e \to I_0^n$ par une application modérée $X \to I_0^n$.
\end{itemize}

5°) est un théorème de prolongement de Tietze, 6°) dit que les fermés modérés
sont des lieux de zéros d'applications modérées ; dans le cadre o-minimal, les
deux sont vrais.\footnote{Voir le livre de van den Dries cité plus haut,
chapitre~8. Rapprochement de l'édition.}

\textbf{La démonstration (page 61).} Si $X$ est un rétracte de $Y$ et de $Z$, de
rétractions $r$ et $r'$, on plonge $Y$ et $Z$ dans le produit fibré $Y
\times_X Z$ (pris au-dessus des rétractions) par $y \mapsto (y, r(y))$ et
$z \mapsto (r'(z), z)$ ; leurs images se coupent exactement le long de la
diagonale de $X$, et l'on prend pour $S'$ ce produit fibré.\footnote{La
transcription marque ce « $Y \times_X Z$ » comme suspect, attendant une somme
amalgamée. Pris au-dessus des rétractions, le produit fibré est juste : les deux
plongements tombent dedans et ne se rencontrent que sur $X$. Le reste de la
ligne est illisible ; on complète l'argument.} On se ramène à ce cas : par
6°), $X$ est l'image réciproque d'un point par $\Psi : Y \to I_0^p$ ; par M5)
et 5°), on plonge $Y$ dans $I_0^n$ par un $\Phi$ prolongeant un plongement
$\varphi$ de $X$ ; on fait de même pour $Z$ ; puis on prolonge $\varphi$ à $Z$
et $\varphi'$ à $Y$. On obtient un cube commutatif de monomorphismes à faces
cartésiennes, où $X$, $Y$, $Z$ se plongent dans des cubes $I^\nu$, $I^N$,
$I^{N'}$ ($\nu = n + n'$, $N = p + \nu$, $N' = p' + \nu$), et $I^\nu$ est un
rétracte de $I^N$ et de $I^{N'}$ ; le cas du rétracte donne $S$.

\textbf{Les fonctions modérées (page 62).} Les conditions 5°) et 6°) ne
s'expriment pas simplement en termes des $\mathcal{M}_n$. Il propose de prendre
pour donnée primitive les ensembles $F_n = \mathrm{Hom}_{\mathcal{M}}(I_0^n,
I_0)$ des fonctions modérées de $n$ variables, avec (F1) les projections sont
dans $F_n$, (F2) les $F_n$ sont stables par composition — c'est la donnée d'une
\emph{théorie de Lawvere} sur une sorte, ou d'un clone d'opérations —, et l'on
pose $\mathcal{M}_n$ = intersections finies de lieux $\varphi^{-1}(\alpha)$,
$\varphi \in F_n$, $\alpha$ un point modéré. Viennent (F3), inachevé, pour
avoir M3), et (F4) : un point $0$ dont les points modérés sont des lieux de
zéros, la réunion des deux axes dans $\mathcal{M}_2$, deux points distincts.
\footnote{Sous (F4) a), il écrit « $\forall \alpha \in F_1$ » alors que
$\alpha$ est un point, donc dans $F_0$ ; on lit $F_0$.} Une fonction est dans
$F_n$ si et seulement si son graphe est dans $\mathcal{M}_{n+1}$ : le « il
faut » demande que la diagonale de $I \times I$ soit dans $\mathcal{M}_2$, le
« il suffit » doit être postulé.

\pagerange{64}{67}
\section*{Axiomes dans les cubes ; carrés de recollement (pages 64 à 67)}

\textbf{Une liste d'axiomes (page 64).} Avec $I = [-1, 1]$ : les
$\mathcal{M}_n$ sont stables par intersections et réunions finies et par images
directes et réciproques d'applications de coordonnées ; $\mathcal{M}_0$
contient $\{0\}$ et $\{1\}$ ; entre crochets, $I$ est injectif et toute partie
modérée est un lieu de zéros $\bigcap_i f_i^{-1}(0)$ de fonctions « permises ».
S'y ajoute une exigence de \emph{triangulation} : la réalisation géométrique de
tout complexe simplicial fini porte une structure modérée admissible, tout objet
est triangulable, et pour $Y \subset X$ une triangulation de $Y$ est induite
par une triangulation de $X$. Exemple : pour des points modérés $a_0, \dots,
a_n$, la partie $\{ (x, y) : x_i \geqslant 0, \sum x_i = 1, y = \sum a_i x_i
\}$ — le graphe de la fonction affine sur le simplexe qui vaut $a_i$ au sommet
$i$ — est dans $\mathcal{M}_{n+2}$.\footnote{La condition a) sur les simplexes
types, ligne surchargée de biffures, ne se lit que par fragments. Le
théorème de triangulation des ensembles définissables, compatible avec une
famille finie de parties, est un des théorèmes centraux de la théorie
o-minimale.}

\textbf{Quand un carré est-il une somme amalgamée ? (page 65).} Soit un carré
commutatif $X \overset{i}{\hookrightarrow} Y$, $X \xrightarrow{p} Z$, $Z
\overset{i'}{\hookrightarrow} S$, $Y \xrightarrow{p'} S$. Il énonce :

\begin{itemize}
\item[a)] $i'$ est un monomorphisme ;
\item[b)] le carré est cartésien : $X = p'^{-1}(Z)$ ;
\item[c)] $Y \times_S Y = \Delta_Y \cup (X \times_Z X)$ ;
\item[d)] $Z \sqcup Y \to S$ est un épimorphisme effectif ;
\end{itemize}

entraînent $S \simeq Y \sqcup_X Z$. C'est exact : la paire noyau de
$Z \sqcup Y \to S$ se décompose, par a) et b), en $Z \sqcup X \sqcup X \sqcup
(Y \times_S Y)$, et par c) une application $(f, g) : Z \sqcup Y \to T$ avec
$f \circ p = g \circ i$ égalise cette paire, donc passe au quotient par
d).\footnote{L'argument est celui de la page, écrit en formules. Dans c), la
réunion n'est pas disjointe — les deux termes se rencontrent le long de
$\Delta_X$ — alors que la page écrit une fois $(X \times_Z X) \sqcup \Delta_Y$.
Le calcul de la paire noyau suppose les sommes disjointes et universelles.}

\textbf{Le réaliser dans un cube (pages 66 et 67).} Il cherche à réaliser un tel
carré avec $S$ un cube, en plongeant $Z$ et $Y$ par des fonctions $f_\alpha$,
$g_\alpha$ compatibles ($f_\alpha \circ p = g_\alpha \circ i$, prolongées par
injectivité) et en ajoutant des fonctions qui s'annulent exactement sur $X$
(par 6°)). Il note « OK pour a) », « OK pour b) ». Pour c), il passe par une
factorisation $S = I^{p+q} \to S' = I^q$ et écrit $Y \times_{I^q} Y = \Delta_Y
\cup (X \times X)$ ; mais cette égalité demande que $Y \to I^q$ soit injective
hors de $X$, ce que la construction n'assure pas, et la page écrit elle-même
un « $\neq$ » pour l'égalité analogue sur le cube.\footnote{Ce « $\neq$ » est
d'une lecture douteuse. Les lettres de $f_\alpha$, $g_\alpha$ sont
surchargées.} La condition c) reste ouverte, et avec elle l'existence générale
des sommes amalgamées.

La moitié gauche de la page 65, écrite en travers, est sans rapport : la
présentation de la cohomologie d'un fibré en grassmanniennes par les classes de
Chern du sous-fibré et du quotient tautologiques, $H(S)[c_i(M), c_j(N)]$ modulo
$c(M)\,c(N) = c(E)$, et la formule correspondante sur un fibré en
drapeaux.\footnote{Avec des lieux $D_\sigma(E)$ et $\Sigma_{n,\sigma}(E)$ qui
évoquent des lieux de dégénérescence de type Schubert ; un exposant de $c_i$ et
un indice de sommation sont illisibles.}

\pagerange{69}{72}
\section*{Les $C$-modèles (pages 69 à 72)}

La dernière mouture, mise au propre à l'encre noire, numérotée « 1) » et
« 2) », reprend tout sur un ensemble $R$ (« ce sera bientôt $\mathbb{R}$ »).

\subsection*{1. Les classes $C_n$ (pages 69 et 70)}

Pour chaque $n$, soit $C_n$ un ensemble de parties de $R^n$,\footnote{La page
écrit « $C_n \in R^n$ » ; on lit $C_n \subset \mathfrak{P}(R^n)$, comme la
suite de la page.} avec

\begin{itemize}
\item[(1)] $C_n$ est stable par permutation des coordonnées ;
\item[(2)] l'image d'un élément de $C_n$ par une projection $R^n \to R^p$ est
  dans $C_p$ ;
\item[(3)] l'image d'un élément de $C_n$ par une application qui répète les
  coordonnées, $(x_1, \dots, x_n) \mapsto (x_1, \dots, x_1, \dots, x_n, \dots,
  x_n)$, est dans $C_{p_1 + \cdots + p_n}$ ;
\item[(4)] si $X \in C_m$ et $Y \in C_n$, $X \times Y \in C_{m+n}$.
\end{itemize}

De façon équivalente, on se donne pour chaque ensemble fini $I$ une classe
$C_I$ de parties de $R^I$, stable par les applications $\alpha^* : R^J \to
R^I$ induites par les applications d'ensembles finis $\alpha : I \to J$, et par
produits.

\textbf{Proposition.} \emph{Si $X, Y \in C_n$, alors $X \cap Y \in C_n$} —
pourvu que les $C_I$ soient aussi stables par \emph{images réciproques} par les
$\alpha^*$ : alors $X \cap Y$ est l'image réciproque de $X \times Y$ par la
diagonale $R^n \to R^{2n}$.\footnote{Hypothèse ajoutée. Sous (1) à (4)
seuls, l'énoncé est faux : les classes engendrées par les deux parties
$[0,2]$ et $[1,3]$ de $\mathbb{R}$ ne contiennent dans $C_1$ que ces deux
parties, car (1) à (4) ne font que permuter, répéter et oublier des
coordonnées de produits ; $[1,2]$ n'y est pas. La stabilité par images
réciproques est l'axiome M$'$2) de la page 55, que cette mouture a omis.}

Une application $f : X \to Y$, $X \in C_I$, $Y \in C_J$, est un
\emph{$C$-morphisme} si son graphe est dans $C_{I \sqcup J}$.\footnote{Il
accole à $C$, dans « $C$-morphisme » et « $C$-modèle », une petite marque en
indice (étoile ou croix), transcrite $C_*$ ; on écrit simplement $C$.} Les
identités, les composés, les couples $(g, h) : X \to Y \times Z$ dont les
composantes le sont, les restrictions des $\alpha^*$ et les inclusions sont des
$C$-morphismes ; l'image directe et l'image réciproque d'une partie de $C$ par
un $C$-morphisme sont dans $C$ ; et l'inverse d'un $C$-morphisme bijectif est
un $C$-morphisme, son graphe étant le transposé du premier
(1).\footnote{L'énoncé c) de la page, sur l'inverse d'une bijection, finit sur
une condition illisible ; on énonce ce qui est vrai. Les énoncés e) à h) sont
en marge, en biais ; leurs lettres et indices sont peu sûrs. Plusieurs de ces
énoncés — composés, images — passent par des intersections, donc par
l'hypothèse ajoutée ci-dessus.}

\textbf{Les $C$-modèles.} Une structure de $C$-modèle sur un ensemble $E$ est
la donnée, pour chaque $n$, d'un ensemble $T_n$ d'injections $E \to R^n$, non
tous vides, tels que, si $u \in T_n$, une injection $v : E \to R^m$ est dans
$T_m$ si et seulement si $v(E) \in C_m$ et $v \circ u^{-1}$ est un
$C$-morphisme. En marge, une définition équivalente par l'ensemble des
$C$-morphismes $E \to R$. Tout $X \in C_n$ est un $C$-modèle, et les
$C$-modèles forment une catégorie, avec produits finis, où un morphisme est un
isomorphisme si et seulement s'il est bijectif. Il y a un $C$-modèle à un
point si et seulement si un $C_n$ contient une partie non vide, si et seulement
si $R^0 \in C_0$ ; on demande plus fortement

\begin{itemize}
\item[(5)] $R^n$ est réunion des éléments de $C_n$.
\end{itemize}

\subsection*{2. Le cas topologique (pages 71 et 72)}

On suppose $R$ topologique séparé et

\begin{itemize}
\item[(7)] les éléments de $C_n$ sont compacts.
\end{itemize}

Alors tout $C$-morphisme est continu — son graphe est un fermé du produit, à
valeurs dans un compact séparé —, et l'on a un foncteur des $C$-modèles vers
les espaces compacts : c'est la catégorie $\mathcal{M}$ des pages 53 à 57,
retrouvée de l'autre côté. On suppose encore

\begin{itemize}
\item[(8)] pour $X \in C_n$ et $x \in X$, les $Y \in C_n$ qui sont des
  voisinages de $x$ dans $X$ forment un système fondamental de voisinages
\end{itemize}

— il suffit que les intérieurs des éléments de $C_n$ forment une base de
$R^n$ —, et

\begin{itemize}
\item[(6)] $C_n$ est stable par réunions finies (donc contient $\emptyset$).
\end{itemize}

\textbf{Proposition (page 72).} \emph{Si un $C$-modèle $X$ est réunion d'une
famille finie de $C$-parties $X_i$, une application $f : X \to Y$ est un
$C$-morphisme si et seulement si chaque $f|X_i$ en est un} : le graphe de $f$
est la réunion des graphes des restrictions (6), et chacun de ceux-ci est la
trace du graphe de $f$ sur $X_i \times Y$. \emph{Corollaire} : $X$ est la somme
amalgamée des $X_i$ au-dessus des $X_i \cap X_j$.

Avec (8), il en déduit que les $C$-morphismes d'un $C$-modèle dans un autre
sont les sections d'un faisceau, celui de leurs germes — ce qui demande
précisément ce recollement fini, donc (6).\footnote{La page 71 affirme le
caractère faisceautique avant d'introduire (6), à la page 72 ; le numéro 6,
cerclé, avait été réservé. On met l'hypothèse à sa place logique.} La
connaissance du faisceau des $C$-morphismes vers $R$ ne donne
pas, dit-il, la structure de modèle.\footnote{« pas » est de lecture douteuse,
et l'addition interlinéaire qui précède est en partie illisible.} D'où la
notion de \emph{$C$-espace} : un espace localement compact muni d'un faisceau
de fonctions continues à valeurs dans $R$ tel que tout point ait un voisinage
compact qui soit, pour la structure induite, un $C$-modèle — l'analogue des
« espaces définissables » de la théorie o-minimale.

\textbf{La question finale.} Pour deux inclusions de $C$-modèles $X
\hookleftarrow Z \hookrightarrow Y$, la somme amalgamée existe-t-elle ? Sa
réponse : « OK », si l'on suppose que pour $Y \subset X$ dans $C_n$, tout
$C$-morphisme $Y \to R$ se prolonge en un $C$-morphisme $X \to R$ — la
condition 5°) de la page 60, sous forme d'un théorème de Tietze. C'est sur
cette condition que le dossier s'arrête.
\end{document}
